\title{2024 年高考数学 天津卷}
\date{}
\maketitle

\section*{一、选择题}
在每小题给出的四个选项中，只有一项是符合题目要求的。

\begin{question}[points=5]
集合 $A = \{1,2,3,4\}$, $B = \{2,3,4,5\}$, 则 $A \cap B =$ ( )
\begin{choices}
  \item $\{1,2,3,4\}$
  \item $\{2,3,4\}$
  \item $\{2,4\}$
  \item $\{1\}$
\end{choices}
\begin{solution}
\textbf{答案：}$\{2,3,4\}$

\textbf{分析：}根据集合交集的概念直接求解即可.

\textbf{详解：}因为集合 $A = \{1,2,3,4\}$, $B = \{2,3,4,5\}$, 所以 $A \cap B = \{2,3,4\}$. 故选 B.
\end{solution}
\end{question}

\begin{question}[points=5]
设 $a,b \in \mathbb{R}$, 则 "$a^3 = b^3$" 是 "$3^a = 3^b$" 的 ( )
\begin{choices}
  \item 充分不必要条件
  \item 必要不充分条件
  \item 充要条件
  \item 既不充分也不必要条件
\end{choices}
\begin{solution}
\textbf{答案：}充要条件

\textbf{分析：}说明二者与同一个命题等价，再得到二者等价，即是充分必要条件.

\textbf{详解：}根据立方的性质和指数函数的性质，$a^3 = b^3$ 和 $3^a = 3^b$ 都当且仅当 $a = b$, 所以二者互为充要条件. 故选 C.
\end{solution}
\end{question}

\begin{question}[points=5]
下列图中，相关性系数最大的是 ( )

\begin{center}
  \includegraphics[width=.38\linewidth]{question-3-1.jpg}
\end{center}
\begin{choices}
  \item A
  \item B
  \item C
  \item D
\end{choices}
\begin{solution}
\textbf{答案：}A

\textbf{分析：}由点的分布特征可直接判断

\textbf{详解：}观察 4 幅图可知，A 图散点分布比较集中，且大体接近某一条直线，线性回归模型拟合效果比较好，呈现明显的正相关，$r$ 值相比于其他 3 图更接近 1. 故选 A.
\end{solution}
\end{question}

\begin{question}[points=5]
下列函数是偶函数的是 ( )
\begin{choices}
  \item $y = \frac{e^x - x^2}{x^2 + 1}$
  \item $y = \frac{\cos x + x^2}{x^2 + 1}$
  \item $y = \frac{e^x - x}{x + 1}$
  \item $y = \frac{\sin x + 4x}{e^{|x|}}$
\end{choices}
\begin{solution}
\textbf{答案：}$y = \frac{\cos x + x^2}{x^2 + 1}$

\textbf{分析：}根据偶函数的判定方法一一判断即可.

\textbf{详解：}对 A, 设 $f(x)=\frac{e^x-x^2}{x^2+1}$, 定义域为 $\mathbb{R}$, 但 $f(-1)=\frac{e^{-1}-1}{2}$, $f(1)=\frac{e-1}{2}$, 则 $f(-1)\neq f(1)$, 故 A 错误. 对 B, 设 $g(x)=\frac{\cos x+x^2}{x^2+1}$, 定义域为 $\mathbb{R}$, 且 $g(-x)=\frac{\cos(-x)+(-x)^2}{(-x)^2+1}=\frac{\cos x+x^2}{x^2+1}=g(x)$, 则 $g(x)$ 为偶函数, 故 B 正确. 对 C, 设 $h(x)=\frac{e^x-x}{x+1}$, 定义域为 $\{x|x\neq -1\}$, 不关于原点对称, 则 $h(x)$ 不是偶函数, 故 C 错误. 对 D, 设 $\varphi(x)=\frac{\sin x+4x}{e^{|x|}}$, 定义域为 $\mathbb{R}$, 因为 $\varphi(1)=\frac{\sin 1+4}{e}$, $\varphi(-1)=\frac{-\sin 1-4}{e}$, 则 $\varphi(1)\neq\varphi(-1)$, 则 $\varphi(x)$ 不是偶函数, 故 D 错误. 故选 B.
\end{solution}
\end{question}

\begin{question}[points=5]
若 $a = 4.2^{-0.3}$, $b = 4.2^{0.3}$, $c = \log_{4.2}{0.2}$, 则 $a$, $b$, $c$ 的大小关系为 ( )
\begin{choices}
  \item $a > b > c$
  \item $b > a > c$
  \item $c > a > b$
  \item $b > c > a$
\end{choices}
\begin{solution}
\textbf{答案：}$b > a > c$

\textbf{分析：}利用指数函数和对数函数的单调性分析判断即可.

\textbf{详解：}因为 $y=4.2^x$ 在 $\mathbb{R}$ 上递增, 且 $-0.3<0<0.3$, 所以 $0<4.2^{-0.3}<4.2^0<4.2^{0.3}$, 即 $0<a<1<b$. 因为 $y=\log_{4.2}x$ 在 $(0,+\infty)$ 上递增, 且 $0<0.2<1$, 所以 $\log_{4.2}{0.2}<\log_{4.2}{1}=0$, 即 $c<0$. 所以 $b>a>c$. 故选 B.
\end{solution}
\end{question}

\begin{question}[points=5]
若 $m,n$ 为两条不同的直线, $\alpha$ 为一个平面, 则下列结论中正确的是 ( )
\begin{choices}
  \item 若 $m // \alpha$, $n \subset \alpha$, 则 $m // n$
  \item 若 $m // \alpha$, $n // \alpha$, 则 $m // n$
  \item 若 $m // \alpha$, $n \perp \alpha$, 则 $m \perp n$
  \item 若 $m // \alpha$, $n \perp \alpha$, 则 $m$ 与 $n$ 相交
\end{choices}
\begin{solution}
\textbf{答案：}若 $m // \alpha$, $n \perp \alpha$, 则 $m \perp n$

\textbf{分析：}根据线面平行的性质可判断AB的正误，根据线面垂直的性质可判断CD的正误.

\textbf{详解：}对于A, 若 $m // \alpha$, $n \subset \alpha$, 则 $m,n$ 平行或异面, 故A错误. 对于B, 若 $m // \alpha$, $n // \alpha$, 则 $m,n$ 平行或异面或相交, 故B错误. 对于C, $m // \alpha$, $n \perp \alpha$, 过 $m$ 作平面 $\beta$, 使得 $\beta \cap \alpha = s$, 因为 $m \subset \beta$, 故 $m // s$, 而 $s \subset \alpha$, 故 $n \perp s$, 故 $m \perp n$, 故C正确. 对于D, 若 $m // \alpha$, $n \perp \alpha$, 则 $m$ 与 $n$ 相交或异面, 故D错误. 故选C.
\end{solution}
\end{question}

\begin{question}[points=5]
已知函数 $f(x)=\sin^3\left(\omega x + \frac{\pi}{3}\right) (\omega>0)$ 的最小正周期为 $\pi$. 则函数在 $\left[-\frac{\pi}{12},\frac{\pi}{6}\right]$ 的最小值是 ( )
\begin{choices}
  \item $-\frac{\sqrt{3}}{2}$
  \item $-\frac{3}{2}$
  \item $0$
  \item $\frac{\sqrt{3}}{2}$
\end{choices}
\begin{solution}
\textbf{答案：}$-\frac{\sqrt{3}}{2}$

\textbf{分析：}先由诱导公式化简, 结合周期公式求出 $\omega$, 得 $f(x)=-\sin 2x$, 再整体求出 $x\in\left[-\frac{\pi}{12},\frac{\pi}{6}\right]$ 时 $2x$ 的范围, 结合正弦三角函数图象特征即可求解.

\textbf{详解：}$f(x)=\sin^3\left(\omega x+\frac{\pi}{3}\right)=\sin(3\omega x+\pi)=-\sin 3\omega x$, 由 $T=\frac{2\pi}{3\omega}=\pi$ 得 $\omega=\frac{2}{3}$, 即 $f(x)=-\sin 2x$. 当 $x\in\left[-\frac{\pi}{12},\frac{\pi}{6}\right]$ 时, $2x\in\left[-\frac{\pi}{6},\frac{\pi}{3}\right]$, 画出 $f(x)=-\sin 2x$ 图象, 由图可知 $f(x)=-\sin 2x$ 在 $\left[-\frac{\pi}{12},\frac{\pi}{6}\right]$ 上递减, 所以当 $x=\frac{\pi}{6}$ 时, $f(x)_{\min}=-\sin\frac{\pi}{3}=-\frac{\sqrt{3}}{2}$. 故选A.
\end{solution}
\end{question}

\begin{question}[points=5]
双曲线 $\frac{x^2}{a^2}-\frac{y^2}{b^2}=1(a>0,b>0)$ 的左、右焦点分别为 $F_1$、$F_2$. $P$ 是双曲线右支上一点, 且直线 $PF_2$ 的斜率为 $2$. $\triangle PF_1F_2$ 是面积为 $8$ 的直角三角形, 则双曲线的方程为 ( )
\begin{choices}
  \item $\frac{x^2}{8}-\frac{y^2}{2}=1$
  \item $\frac{x^2}{8}-\frac{y^2}{4}=1$
  \item $\frac{x^2}{2}-\frac{y^2}{8}=1$
  \item $\frac{x^2}{4}-\frac{y^2}{8}=1$
\end{choices}
\begin{solution}
\textbf{答案：}$\frac{x^2}{2}-\frac{y^2}{8}=1$

\textbf{分析：}可利用 $\triangle PF_1F_2$ 三边斜率问题与正弦定理, 转化出三边比例, 设 $PF_2=m$, 由面积公式求出 $m$, 由勾股定理得出 $c$, 结合第一定义再求出 $a$.

\textbf{详解：}由题可知, 点 $P$ 必落在第四象限, $\angle F_1PF_2=90^\circ$, 设 $PF_2=m$, $\angle PF_2F_1=\theta_1$, $\angle PF_1F_2=\theta_2$, 由 $k_{PF_2}=\tan\theta_1=2$, 得 $\sin\theta_1=\frac{2}{\sqrt{5}}$. 因为 $\angle F_1PF_2=90^\circ$, 所以 $k_{PF_1}\cdot k_{PF_2}=-1$, 得 $k_{PF_1}=-\frac{1}{2}$, 即 $\tan\theta_2=\frac{1}{2}$, $\sin\theta_2=\frac{1}{\sqrt{5}}$. 由正弦定理可得 $PF_1:PF_2:F_1F_2=\sin\theta_1:\sin\theta_2:\sin90^\circ=2:1:\sqrt{5}$, 则由 $PF_2=m$ 得 $PF_1=2m$, $F_1F_2=2c=\sqrt{5}m$. 由 $S_{\triangle PF_1F_2}=\frac{1}{2}PF_1\cdot PF_2=m^2=8$ 得 $m=2\sqrt{2}$, 则 $PF_2=2\sqrt{2}$, $PF_1=4\sqrt{2}$, $F_1F_2=2c=2\sqrt{10}$, $c=\sqrt{10}$. 由双曲线定义可得 $PF_1-PF_2=2a=2\sqrt{2}$, $a=\sqrt{2}$, $b=\sqrt{c^2-a^2}=\sqrt{8}=2\sqrt{2}$, 所以双曲线方程为 $\frac{x^2}{2}-\frac{y^2}{8}=1$. 故选C.
\end{solution}
\end{question}

\begin{question}[points=5]
一个五面体 $ABC-DEF$. 已知 $AD \parallel BE \parallel CF$, 且两两之间距离为 $1$. 并已知 $AD=1$, $BE=2$, $CF=3$. 则该五面体的体积为 ( )
\begin{choices}
  \item $\frac{\sqrt{3}}{6}$
  \item $\frac{3\sqrt{3}}{4}+\frac{1}{2}$
  \item $\frac{\sqrt{3}}{2}$
  \item $\frac{3\sqrt{3}}{4}-\frac{1}{2}$
\end{choices}
\begin{solution}
\textbf{答案：}$\frac{\sqrt{3}}{2}$

\textbf{分析：}采用补形法, 补成一个棱柱, 求出其直截面, 再利用体积公式即可.

\textbf{详解：}用一个完全相同的五面体 $HIJ-LMN$ (顶点与五面体 $ABC-DEF$ 一一对应) 与该五面体相嵌, 使得 $D,N$; $E,M$; $F,L$ 重合. 因为 $AD \parallel BE \parallel CF$, 且两两之间距离为 $1$, $AD=1$, $BE=2$, $CF=3$, 则形成的新组合体为一个三棱柱, 该三棱柱的直截面 (与侧棱垂直的截面) 为边长为 $1$ 的等边三角形, 侧棱长为 $1+3=2+2=3+1=4$, $V_{ABC-DEF}=\frac{1}{2}V_{ABC-HIJ}=\frac{1}{2}\times\frac{1}{2}\times 1\times 1\times \frac{\sqrt{3}}{2}\times 4=\frac{\sqrt{3}}{2}$. 故选C.
\end{solution}
\end{question}

\section*{二、填空题}
本大题共6小题，每小题5分，共30分. 试题中包含两个空的，答对1个的给3分，全部答对的给5分.

\begin{question}[points=5]
已知 $i$ 是虚数单位，复数 $(\sqrt{5}+i)(\sqrt{5}-2i)=$ ．
\fillin[{$7-\sqrt{5}i$}]
\begin{solution}
\textbf{答案：}$7-\sqrt{5}i$

\textbf{分析：}直接利用复数乘法法则计算。

\textbf{详解：}$(\sqrt{5}+i)(\sqrt{5}-2i)=5-2\sqrt{5}i+\sqrt{5}i-2i^2=7-\sqrt{5}i$。
\end{solution}
\end{question}

\begin{question}[points=5]
在 $\left(\frac{3}{x^3}+\frac{x^3}{3}\right)^6$ 的展开式中，常数项为 ．
\fillin[{$20$}]
\begin{solution}
\textbf{答案：}$20$

\textbf{分析：}写出二项展开式的通项，令 $x$ 的指数为零。

\textbf{详解：}通项为 $T_{r+1}=\binom{6}{r}\left(\frac{3}{x^3}\right)^{6-r}\left(\frac{x^3}{3}\right)^r=\binom{6}{r}3^{6-2r}x^{6r-18}$。令 $6r-18=0$，得 $r=3$，所以常数项为 $\binom{6}{3}=20$。
\end{solution}
\end{question}

\begin{question}[points=5]
$(x-1)^2+y^2=25$ 的圆心与抛物线 $y^2=2px(p>0)$ 的焦点 $F$ 重合, $A$ 为两曲线的交点, 则原点到直线 $AF$ 的距离为.
\fillin[{$\frac{4}{5}$}]
\begin{solution}
\textbf{答案：}$\frac{4}{5}$

\textbf{分析：}先求出圆心坐标，从而可求焦准距，再联立圆和抛物线方程，求 $A$ 及 $AF$ 的方程，从而可求原点到直线 $AF$ 的距离.

\textbf{详解：}圆 $(x-1)^2+y^2=25$ 的圆心为 $F(1,0)$, 故 $\frac{p}{2}=1$ 即 $p=2$. 由 $\begin{cases} (x-1)^2+y^2=25 \\ y^2=4x \end{cases}$ 得 $x^2+2x-24=0$, 故 $x=4$ 或 $x=-6$ (舍), 故 $A(4,\pm4)$, 所以直线 $AF$ 的方程为 $y=\pm\frac{4}{3}(x-1)$, 即 $4x-3y-4=0$ 或 $4x+3y-4=0$, 故原点到直线 $AF$ 的距离为 $d=\frac{|4|}{\sqrt{4^2+3^2}}=\frac{4}{5}$.
\end{solution}
\end{question}

\begin{question}[points=5]
$A,B,C,D,E$ 五种活动, 甲、乙都要选择三个活动参加. (1) 甲选到 $A$ 的概率为; 已知乙选了 $A$ 活动, 他再选择 $B$ 活动的概率为.
\fillin[{$\frac{3}{5}$, $\frac{1}{2}$}]
\begin{solution}
\textbf{答案：}$\frac{3}{5}$, $\frac{1}{2}$

\textbf{分析：}结合列举法或组合公式和概率公式可求甲选到 $A$ 的概率; 采用列举法或者条件概率公式可求乙选了 $A$ 活动, 他再选择 $B$ 活动的概率.

\textbf{详解：}解法一: 从五个活动中选三个的情况有 $C_5^3=10$ 种, 其中甲选到 $A$ 有 $C_4^2=6$ 种, 概率为 $\frac{6}{10}=\frac{3}{5}$. 乙选 $A$ 有 $C_4^2=6$ 种, 其中再选 $B$ 有 $C_3^1=3$ 种, 条件概率为 $\frac{3}{6}=\frac{1}{2}$. 解法二: $P(\text{甲选}A)=\frac{C_4^2}{C_5^3}=\frac{3}{5}$, $P(\text{乙选}B|\text{乙选}A)=\frac{C_3^1}{C_4^2}=\frac{1}{2}$.
\end{solution}
\end{question}

\begin{question}[points=5]
在边长为 $1$ 的正方形 $ABCD$ 中, 点 $E$ 为线段 $CD$ 的三等分点, $CE = \frac{1}{2}DE$, $\overrightarrow{BE} = \lambda \overrightarrow{BA} + \mu \overrightarrow{BC}$, 则 $\lambda+\mu =$ ; 若 $F$ 为线段 $BE$ 上的动点, $G$ 为 $AF$ 中点, 则 $\overrightarrow{AF} \cdot \overrightarrow{DG}$ 的最小值为.
\fillin[{$\frac{4}{3}$, $-\frac{5}{18}$}]
\begin{solution}
\textbf{答案：}$\frac{4}{3}$, $-\frac{5}{18}$

\textbf{分析：}解法一: 以 $\overrightarrow{BA},\overrightarrow{BC}$ 为基底, 根据向量的线性运算求 $\overrightarrow{BE}$, 可得 $\lambda+\mu$, 设 $\overrightarrow{BF}=k\overrightarrow{BE}$, 求 $\overrightarrow{AF},\overrightarrow{DG}$, 结合数量积的运算律求最小值. 解法二: 建系标点, 根据向量的坐标运算求 $\overrightarrow{BE}$, 可得 $\lambda+\mu$, 设 $F(a,-3a),a\in[-\frac{1}{3},0]$, 求 $\overrightarrow{AF},\overrightarrow{DG}$, 结合数量积的坐标运算求最小值.

\textbf{详解：}解法一: 因为 $CE=\frac{1}{2}DE$, 即 $\overrightarrow{CE}=\frac{1}{3}\overrightarrow{BA}$, 则 $\overrightarrow{BE}=\overrightarrow{BC}+\overrightarrow{CE}=\frac{1}{3}\overrightarrow{BA}+\overrightarrow{BC}$, 所以 $\lambda=\frac{1}{3},\mu=1$, $\lambda+\mu=\frac{4}{3}$. 由题意 $|\overrightarrow{BC}|=|\overrightarrow{BA}|=1$, $\overrightarrow{BA}\cdot\overrightarrow{BC}=0$. 设 $\overrightarrow{BF}=k\overrightarrow{BE}= \frac{k}{3}\overrightarrow{BA}+k\overrightarrow{BC}$, $k\in[0,1]$. 则 $\overrightarrow{AF}=\overrightarrow{AB}+\overrightarrow{BF}=(\frac{k}{3}-1)\overrightarrow{BA}+k\overrightarrow{BC}$. 又 $G$ 为 $AF$ 中点, $\overrightarrow{DG}=\overrightarrow{DA}+\overrightarrow{AG}=-\overrightarrow{BC}+\frac{1}{2}\overrightarrow{AF}=\frac{1}{2}(\frac{k}{3}-1)\overrightarrow{BA}+(\frac{k}{2}-1)\overrightarrow{BC}$. 则 $\overrightarrow{AF}\cdot\overrightarrow{DG}=[ (\frac{k}{3}-1)\overrightarrow{BA}+k\overrightarrow{BC} ]\cdot[\frac{1}{2}(\frac{k}{3}-1)\overrightarrow{BA}+(\frac{k}{2}-1)\overrightarrow{BC}] =\frac{1}{2}(\frac{k}{3}-1)^2 + k(\frac{k}{2}-1)=\frac{5}{9}(k-\frac{6}{5})^2-\frac{3}{10}$, 又 $k\in[0,1]$, 当 $k=1$ 时取最小值 $-\frac{5}{18}$.
\end{solution}
\end{question}

\begin{question}[points=5]
若函数 $f(x)=\sqrt{2x^2-ax} - |ax-2| + 1$ 有唯一零点, 则 $a$ 的取值范围为.
\fillin[{$(-\sqrt{3},-1)\cup(1,\sqrt{3})$}]
\begin{solution}
\textbf{答案：}$(-\sqrt{3},-1)\cup(1,\sqrt{3})$

\textbf{分析：}结合函数零点与两函数的交点的关系, 构造函数 $g(x)=\sqrt{2x^2-ax}$ 与 $h(x)=\begin{cases} ax-3, & x\geq\frac{2}{a} \\ 1-ax, & x<\frac{2}{a} \end{cases}$, 则两函数图象有唯一交点, 分 $a=0$、$a>0$ 与 $a<0$ 进行讨论.

\textbf{详解：}令 $f(x)=0$, 即 $\sqrt{2x^2-ax}=|ax-2|-1$. 由题 $x^2-ax\geq0$. 当 $a=0$ 时, $\sqrt{2x^2}=1$, $x=\pm\frac{\sqrt{2}}{2}$, 两解, 舍去. 当 $a>0$ 时, 则 $\sqrt{2x^2-ax}=|ax-2|-1=\begin{cases} ax-3, & x\geq\frac{2}{a}\\ 1-ax, & x<\frac{2}{a} \end{cases}$. 由 $x^2-ax\geq0$ 得 $x\geq a$ 或 $x\leq0$. 当 $x\leq0$ 时, $|ax-2|=2-ax$, 方程化为 $\sqrt{2x^2-ax}=1-ax$, 平方得 $(4-a^2)x^2-2ax-1=0$, 即 $[(2+a)x+1][(2-a)x-1]=0$. 当 $a=2$ 时, $4x+1=0$, $x=-\frac{1}{4}$; 当 $0<a<2$ 时, $x=-\frac{1}{2+a}$ 或 $x=\frac{1}{2-a}>0$ (舍); 当 $a>2$ 时, $x=-\frac{1}{2+a}<0$ 或 $x=\frac{1}{2-a}<0$, 两解, 舍去. 所以当 $0<a\leq2$ 时, 在 $x\leq0$ 时有唯一解. 则当 $0<a\leq2$ 时, 需在 $x\geq a$ 时无解. 当 $x\geq a$ 时, $h(x)=ax-3$, $g(x)=\sqrt{2x^2-ax}$, 令 $g(x)=0$ 得 $x=a$ (舍0). $g(x)$ 在 $(a,+\infty)$ 上单调递增. 要无交点, 则 $\frac{1}{a}<a$ 且 $\frac{3}{a}>a$, 解得 $1<a<\sqrt{3}$. 故 $1<a<\sqrt{3}$ 符合. 当 $a<0$ 时, 同理得 $-\sqrt{3}<a<-1$. 综上, $a\in(-\sqrt{3},-1)\cup(1,\sqrt{3})$.
\end{solution}
\end{question}

\section*{三、解答题}
本大题共5小题，共75分. 解答应写出文字说明，证明过程或演算步骤.

\begin{question}[points=14]
在 $\triangle ABC$ 中, $\cos B = \frac{9}{16}$, $b=5$, $\frac{a}{c}=\frac{2}{3}$.

(1) 求 $a$;

(2) 求 $\sin A$;

(3) 求 $\cos(B-2A)$.
\begin{solution}
\textbf{答案：}(1) $a=4$; (2) $\sin A = \frac{\sqrt{7}}{4}$; (3) $\frac{57}{64}$

\textbf{分析：}(1) 设 $a=2t,c=3t$, 利用余弦定理即可得到方程, 解出即可. (2) 法一: 求出 $\sin B$, 再利用正弦定理; 法二: 利用余弦定理求出 $\cos A$, 则得到 $\sin A$. (3) 法一: 根据大边对大角确定 $A$ 为锐角, 则得到 $\cos A$, 再利用二倍角公式和两角差的余弦公式即可. 法二: 直接利用二倍角公式和两角差的余弦公式即可.

\textbf{详解：}(1) 设 $a=2t$, $c=3t$, $t>0$, 则根据余弦定理得 $b^2=a^2+c^2-2ac\cos B$, 即 $25=4t^2+9t^2-2\times2t\times3t\times\frac{9}{16}$, 解得 $t=2$ (负舍), 则 $a=4$, $c=6$.
(2) 法一: 因为 $B$ 为三角形内角, 所以 $\sin B = \sqrt{1-\cos^2 B} = \sqrt{1-\left(\frac{9}{16}\right)^2} = \frac{5\sqrt{7}}{16}$, 再根据正弦定理 $\frac{a}{\sin A}=\frac{b}{\sin B}$, 即 $\frac{4}{\sin A}=\frac{5}{\frac{5\sqrt{7}}{16}}$, 解得 $\sin A = \frac{\sqrt{7}}{4}$.
(3) 法一: 因为 $\cos B=\frac{9}{16}>0$, 且 $B\in(0,\pi)$, 所以 $B\in(0,\frac{\pi}{2})$. 由(2)法一知 $\sin B=\frac{5\sqrt{7}}{16}$. 因为 $a<b$, 则 $A<B$, 所以 $\cos A=\sqrt{1-\sin^2 A}=\sqrt{1-\left(\frac{\sqrt{7}}{4}\right)^2}=\frac{3}{4}$. 则 $\sin2A=2\sin A\cos A=2\times\frac{\sqrt{7}}{4}\times\frac{3}{4}=\frac{3\sqrt{7}}{8}$, $\cos2A=2\cos^2 A-1=2\times\left(\frac{3}{4}\right)^2-1=\frac{1}{8}$. $\cos(B-2A)=\cos B\cos2A+\sin B\sin2A=\frac{9}{16}\times\frac{1}{8}+\frac{5\sqrt{7}}{16}\times\frac{3\sqrt{7}}{8}=\frac{9}{128}+\frac{105}{128}=\frac{114}{128}=\frac{57}{64}$.
\end{solution}
\end{question}

\begin{question}[points=15]
已知四棱柱 $ABCD-A_1B_1C_1D_1$ 中, 底面 $ABCD$ 为梯形, $AB//CD$, $A_1A\perp$ 平面 $ABCD$, $AD\perp AB$, 其中 $AB=AA_1=2$, $AD=DC=1$. $N$ 是 $B_1C_1$ 的中点, $M$ 是 $DD_1$ 的中点.

(1) 求证 $D_1N//$ 平面 $CB_1M$;

(2) 求平面 $CB_1M$ 与平面 $BB_1CC_1$ 的夹角余弦值;

(3) 求点 $B$ 到平面 $CB_1M$ 的距离.
\begin{solution}
\textbf{答案：}(1) 证明见解析; (2) $\frac{2\sqrt{22}}{11}$; (3) $\frac{2\sqrt{11}}{11}$

\textbf{分析：}(1) 取 $CB_1$ 中点 $P$, 连接 $NP$, $MP$, 借助中位线的性质与平行四边形性质定理可得 $D_1N//MP$, 结合线面平行判定定理即可得证. (2) 建立空间直角坐标系, 计算两平面的法向量, 再计算夹角余弦. (3) 利用点到平面的距离公式.

\textbf{详解：}(1) 取 $CB_1$ 中点 $P$, 连接 $NP$, $MP$. 由 $N$ 是 $B_1C_1$ 的中点, 故 $NP//CC_1$, 且 $NP=\frac{1}{2}CC_1$. 由 $M$ 是 $DD_1$ 的中点, 故 $D_1M=\frac{1}{2}DD_1=\frac{1}{2}CC_1$, 且 $D_1M//CC_1$, 则有 $D_1M//NP$, $D_1M=NP$, 故四边形 $D_1MPN$ 是平行四边形, 故 $D_1N//MP$. 又 $MP\subset$ 平面 $CB_1M$, $D_1N\not\subset$ 平面 $CB_1M$, 故 $D_1N//$ 平面 $CB_1M$.
(2) 以 $A$ 为原点建立空间直角坐标系, 有 $A(0,0,0)$, $B(2,0,0)$, $B_1(2,0,2)$, $M(0,1,1)$, $C(1,1,0)$, $C_1(1,1,2)$. 则 $\overrightarrow{CB_1}=(1,-1,2)$, $\overrightarrow{CM}=(-1,0,1)$, $\overrightarrow{BB_1}=(0,0,2)$. 设平面 $CB_1M$ 的法向量 $\vec{m}=(x_1,y_1,z_1)$, 由 $\begin{cases} \vec{m}\cdot\overrightarrow{CB_1}=x_1-y_1+2z_1=0\\ \vec{m}\cdot\overrightarrow{CM}=-x_1+z_1=0 \end{cases}$, 取 $x_1=1$, 则 $z_1=1$, $y_1=3$, 即 $\vec{m}=(1,3,1)$. 设平面 $BB_1CC_1$ 的法向量 $\vec{n}=(x_2,y_2,z_2)$, 由 $\begin{cases} \vec{n}\cdot\overrightarrow{CB_1}=x_2-y_2+2z_2=0\\ \vec{n}\cdot\overrightarrow{BB_1}=2z_2=0 \end{cases}$, 取 $x_2=1$, 则 $y_2=1$, $z_2=0$, 即 $\vec{n}=(1,1,0)$. 则 $\cos\langle\vec{m},\vec{n}\rangle=\frac{\vec{m}\cdot\vec{n}}{|\vec{m}||\vec{n}|}=\frac{1+3}{\sqrt{11}\cdot\sqrt{2}}=\frac{4}{\sqrt{22}}=\frac{2\sqrt{22}}{11}$. 故平面 $CB_1M$ 与平面 $BB_1CC_1$ 的夹角余弦值为 $\frac{2\sqrt{22}}{11}$.
(3) 由 $\overrightarrow{BB_1}=(0,0,2)$, 平面 $CB_1M$ 的法向量 $\vec{m}=(1,3,1)$, 则点 $B$ 到平面 $CB_1M$ 的距离 $d=\frac{|\overrightarrow{BB_1}\cdot\vec{m}|}{|\vec{m}|}=\frac{|2|}{\sqrt{11}}=\frac{2\sqrt{11}}{11}$.
\end{solution}
\end{question}

\begin{question}[points=15]
已知椭圆 $\frac{x^2}{a^2}+\frac{y^2}{b^2}=1(a>b>0)$ 的离心率 $e=\frac{1}{2}$. 左顶点为 $A$, 下顶点为 $B$, $C$ 是线段 $OB$ 的中点, 其中 $S_{\triangle ABC}=\frac{3\sqrt{3}}{2}$.

(1) 求椭圆方程.

(2) 过点 $(0,-\frac{3}{2})$ 的动直线与椭圆有两个交点 $P$, $Q$. 在 $y$ 轴上是否存在点 $T$ 使得 $\overrightarrow{TP}\cdot\overrightarrow{TQ}\leq 0$ 恒成立. 若存在求出这个 $T$ 点纵坐标的取值范围, 若不存在请说明理由.
\begin{solution}
\textbf{答案：}(1) $\frac{x^2}{12}+\frac{y^2}{9}=1$; (2) 存在 $T(0,t)$ 且 $-3\leq t\leq \frac{3}{2}$, 使得 $\overrightarrow{TP}\cdot\overrightarrow{TQ}\leq 0$ 恒成立.

\textbf{分析：}(1) 根据椭圆的离心率和三角形的面积可求基本量, 从而可得椭圆的标准方程. (2) 设直线方程, 与椭圆联立, 利用韦达定理和向量数量积的坐标运算, 结合恒成立条件求 $t$ 的范围.

\textbf{详解：}(1) 因为 $e=\frac{1}{2}$, 故 $a=2c$, $b=\sqrt{3}c$. 所以 $A(-2c,0)$, $B(0,-\sqrt{3}c)$, $C(0,-\frac{\sqrt{3}}{2}c)$. $S_{\triangle ABC}=\frac{1}{2}\times 2c\times \frac{\sqrt{3}}{2}c=\frac{\sqrt{3}}{2}c^2=\frac{3\sqrt{3}}{2}$, 解得 $c=\sqrt{3}$, 所以 $a=2\sqrt{3}$, $b=3$, 故椭圆方程为 $\frac{x^2}{12}+\frac{y^2}{9}=1$.
(2) 设过点 $(0,-\frac{3}{2})$ 的动直线方程为 $y=kx-\frac{3}{2}$, $P(x_1,y_1)$, $Q(x_2,y_2)$, $T(0,t)$. 联立 $\begin{cases} 3x^2+4y^2=36 \\ y=kx-\frac{3}{2} \end{cases}$, 得 $(3+4k^2)x^2-12kx-27=0$, $\Delta>0$, $x_1+x_2=\frac{12k}{3+4k^2}$, $x_1x_2=-\frac{27}{3+4k^2}$. $\overrightarrow{TP}=(x_1,y_1-t)$, $\overrightarrow{TQ}=(x_2,y_2-t)$. $\overrightarrow{TP}\cdot\overrightarrow{TQ}=x_1x_2+(y_1-t)(y_2-t)=x_1x_2+(kx_1-\frac{3}{2}-t)(kx_2-\frac{3}{2}-t)=(1+k^2)x_1x_2-k(\frac{3}{2}+t)(x_1+x_2)+(\frac{3}{2}+t)^2$. 代入韦达定理, 得 $\frac{[(3+2t)^2-12t-45]k^2+3(\frac{3}{2}+t)^2-27}{3+4k^2}$. 要使 $\overrightarrow{TP}\cdot\overrightarrow{TQ}\leq 0$ 恒成立, 需 $\begin{cases} (3+2t)^2-12t-45\leq 0 \\ 3(\frac{3}{2}+t)^2-27\leq 0 \end{cases}$, 解得 $-3\leq t\leq \frac{3}{2}$. 当斜率不存在时, 直线为 $x=0$, 与椭圆交于 $(0,3)$, $(0,-3)$, 此时需 $-3\leq t\leq 3$, 结合得 $-3\leq t\leq \frac{3}{2}$. 综上, 存在 $T(0,t)$ 且 $-3\leq t\leq \frac{3}{2}$, 使得 $\overrightarrow{TP}\cdot\overrightarrow{TQ}\leq 0$ 恒成立.
\end{solution}
\end{question}

\begin{question}[points=15]
已知数列 $\{a_n\}$ 是公比大于 $0$ 的等比数列. 其前 $n$ 项和为 $S_n$. 若 $a_1=1$, $S_2=a_3-1$.

(1) 求数列 $\{a_n\}$ 前 $n$ 项和 $S_n$;

(2) 设 $b_n=\begin{cases} k, & n=a_k \\ b_{n-1}+2k, & a_k<n<a_{k+1} \end{cases}$, $b_1=1$, 其中 $k$ 是大于 $1$ 的正整数.

(i) 当 $n=a_{k+1}$ 时, 求证: $b_{n-1}\geq a_k\cdot b_n$;

(ii) 求 $\sum_{i=1}^{S_n} b_i$.
\begin{solution}
\textbf{答案：}(1) $S_n=2^n-1$; (2) (i) 证明见解析; (ii) $\sum_{i=1}^{S_n} b_i = \frac{(3n-1)4^n+1}{9}$

\textbf{分析：}(1) 设公比为 $q$, 利用 $S_2=a_3-1$ 求出 $q$, 再求和公式. (2) (i) 由 $a_k=2^{k-1}$, $n=2^k$, 推导 $b_{n-1}$ 和 $b_n$, 利用作差法证明不等式. (ii) 分段求和, 利用等差数列求和公式和裂项相消法.

\textbf{详解：}(1) 设等比数列 $\{a_n\}$ 的公比为 $q>0$, 由 $a_1=1$, $S_2=a_3-1$ 得 $1+q=q^2-1$, 即 $q^2-q-2=0$, 解得 $q=2$ (舍负). 所以 $a_n=2^{n-1}$, $S_n=\frac{1-2^n}{1-2}=2^n-1$.
(2) (i) 由 (1) 知 $a_k=2^{k-1}$, $a_{k+1}=2^k$. 当 $n=a_{k+1}=2^k$ 时, $a_k< n-1 < a_{k+1}$, 故 $b_n=b_{2^k}=k+1$, 而 $b_{n-1}=b_{2^k-1}=b_{a_k}+(2^k-1-2^{k-1})\cdot 2k = k + (2^{k-1}-1)\cdot 2k = k + 2^k k - 2k = k(2^k-1)$. 所以 $b_{n-1}-a_kb_n = k(2^k-1) - 2^{k-1}(k+1) = 2^{k-1}k - k - 2^{k-1} = 2^{k-1}(k-1) - k \geq 2(k-1)-k = k-2 \geq 0$ (当且仅当 $k=2$ 等号成立). 故 $b_{n-1}\geq a_k b_n$.
(ii) 由 $S_n=2^n-1$. 求和时, 考虑分段. 当 $2^{k-1}\leq i \leq 2^k-1$ 时, $b_i$ 构成等差数列, 首项 $b_{2^{k-1}}=k$, 公差 $2k$, 项数 $2^{k-1}$. 则 $\sum_{i=2^{k-1}}^{2^k-1} b_i = k\cdot 2^{k-1} + \frac{2^{k-1}(2^{k-1}-1)}{2}\cdot 2k = k\cdot 2^{k-1} + k\cdot 2^{k-1}(2^{k-1}-1) = k\cdot 2^{2k-2} = \frac{1}{9}[(3k-1)4^k - (3k-4)4^{k-1}]$ (经检验成立). 则 $\sum_{i=1}^{S_n} b_i = 1 + \sum_{k=2}^n \sum_{i=2^{k-1}}^{2^k-1} b_i = 1 + \frac{1}{9}\sum_{k=2}^n [(3k-1)4^k - (3k-4)4^{k-1}] = 1 + \frac{1}{9}[(3n-1)4^n - (3\cdot2-4)4^1] = \frac{(3n-1)4^n+1}{9}$. 当 $n=1$ 时, $S_1=1$, $\sum b_i=1$, 公式也成立. 故 $\sum_{i=1}^{S_n} b_i = \frac{(3n-1)4^n+1}{9}$.
\end{solution}
\end{question}

\begin{question}[points=16]
设函数 $f(x)=x\ln x$.

(1) 求 $f(x)$ 图象上点 $(1,f(1))$ 处的切线方程;

(2) 若 $f(x)\geq a(x-\sqrt{x})$ 在 $x\in(0,+\infty)$ 时恒成立, 求 $a$ 的取值范围;

(3) 若 $x_1,x_2\in(0,1)$, 证明 $|f(x_1)-f(x_2)|\leq \sqrt{|x_1-x_2|}$.
\begin{solution}
\textbf{答案：}(1) $y=x-1$; (2) $\{2\}$; (3) 证明见解析

\textbf{分析：}(1) 直接使用导数的几何意义. (2) 先由题设条件得到 $a=2$, 再证明 $a=2$ 时条件满足. (3) 先确定 $f(x)$ 的单调性, 再对 $x_1,x_2$ 分类讨论.

\textbf{详解：}(1) $f'(x)=\ln x+1$, $f(1)=0$, $f'(1)=1$, 切线方程为 $y=x-1$.
(2) 设 $g(t)=a(t-1)-2\ln t$, 则 $f(x)-a(x-\sqrt{x})=x\left[a\left(\frac{1}{\sqrt{x}}-1\right)-2\ln\frac{1}{\sqrt{x}}\right]=x\cdot g\left(\frac{1}{\sqrt{x}}\right)$. 故恒成立等价于 $g(t)\geq 0$ 对 $t>0$ 恒成立. 取 $t=2$ 得 $a\geq 1$, 取 $t=\frac{2}{a}$ 得 $a=2$. 反之, 当 $a=2$ 时, $g(t)=2(t-1)-2\ln t=2(t-1-\ln t)\geq 0$ (因为 $t-1\geq\ln t$). 故 $a$ 的取值范围是 $\{2\}$.
(3) 先证引理: 对 $0<a<b$, $\ln a+1<\frac{f(b)-f(a)}{b-a}<\ln b+1$. 由 $f'(x)=\ln x+1$, $f(x)$ 在 $(0,\frac{1}{e})$ 递减, 在 $(\frac{1}{e},+\infty)$ 递增. 不妨设 $x_1\leq x_2$. 分三种情况: (i) $\frac{1}{e}\leq x_1\leq x_2<1$: $|f(x_1)-f(x_2)|=f(x_2)-f(x_1)<(\ln x_2+1)(x_2-x_1)<x_2-x_1\leq \sqrt{x_2-x_1}$; (ii) $0<x_1\leq x_2\leq \frac{1}{e}$: 利用构造的 $\varphi(x)=x\ln x - c\ln c - \sqrt{c-x}$, 可证 $x_1\ln x_1 - x_2\ln x_2 \leq \sqrt{x_2-x_1}$; (iii) $0<x_1\leq \frac{1}{e}\leq x_2<1$: 利用中间值 $\frac{1}{e}$ 和 (i)(ii) 的结论可得 $|f(x_1)-f(x_2)|\leq \max\{|f(x_1)-f(\frac{1}{e})|, |f(\frac{1}{e})-f(x_2)|\}\leq \sqrt{\frac{1}{e}-x_1}\leq \sqrt{x_2-x_1}$ 或 $\sqrt{x_2-\frac{1}{e}}\leq \sqrt{x_2-x_1}$. 故结论成立.
\end{solution}
\end{question}

